\chapter{总结与展望}
\section*{一、论文工作总结}
随着具身智能技术的飞速发展，机器人正加速从结构化的工业生产线走向非结构化、任务高度并行的居家与社会服务场景。然而，真实物理环境的不确定性、专家示教数据的匮乏以及任务指令的语义模糊性，对机器人模仿学习系统的感知深度与执行自主性提出了严峻挑战。本文针对这些痛点，深入开展了基于观测输入的机器人模仿学习研究，构建了从静态图像技能逆向获取到异构约束实时感知，再到通用多模态决策的完整技术体系。本文的主要工作总结如下：

(1) 针对书法等精密技能示教中第一人称视频稀缺且静态图像丢失动态信息的难题，本文设计并实现了一种基于图像观测输入的书法技能模仿学习框架。该框架通过设计八邻域优先级算法实现了对汉字笔画拓扑骨架的精准分割，并引入结合主动曲线轴机制的模糊高斯混合模型 (FGMM)，利用非线性轴线替代传统的线性主轴，解决了高曲率笔画在时序编码过程中的拟合难题。通过偏差分析 (DA) 算法，系统成功从像素分布中重构出表征笔画粗细变化的协方差序列，并利用高度映射机制将其转化为毛笔的法向控制指令。结合动态运动原语 (DMP) 对轨迹进行建模，实现了书法技能在不同空间尺度下的自适应泛化，确保了机器人在不同书写介质上的几何结构完整性与物理复现保真度。

(2) 为了解决生成式扩散策略在非结构化环境下难以实时遵循复杂环境约束的问题，本文设计了面向异构约束观测输入的模块化残差学习网络 ConstrainNet。该网络通过编码空间避障或运动学限制等异构观测信息，利用零卷积初始化与残差注入机制，在冻结的扩散策略骨干网络中引入引导偏置。这种设计确保了系统在训练初期能够保持预训练的操作先验，避免了全参数微调带来的灾难性遗忘。实验论证了空间约束与运动学约束在特征空间中的解耦特性，实现了多约束条件的零样本组合与即插即用式适配。线性探测与可视化分析表明，该网络成功在特征空间中重构了避障所需的几何语义，使机器人能够从被动的反应式滤波转向主动的全局轨迹规划，显著提升了作业安全性。

(3) 为了提升机器人系统在动态交互环境中的自主感知与决策能力，本文构建了基于多模态观测输入的闭环模仿学习系统。首先，利用自监督学习策略构建了语音去噪网络，通过最小化两组带噪信号间的差异，在无需干净标签的前提下实现了对环境指令的特征重建，并将其作为高层意图开关触发 ConstrainNet 模块调度。其次，通过集成 YOLOv8s 目标检测模型，系统能够直接从原始图像流中实时提取障碍物的四维边界框信息，并将其转化为策略网络的连续约束输入，实现了从被动指令触发向主动闭环感知的演进。在长时序分拣任务实验中，该系统展示了规避动态障碍物的能力，验证了残差注入机制在处理非稳态工况时的实时性与可靠性。

(4) 针对大规模机器人预训练中动作空间异构性的挑战，本文开发了端到端训练的通用视觉-语言-动作 (VLA) 模仿学习系统。系统设计了基于对比学习的潜动作编码器，利用监督对比损失与均匀性损失将不同构型机器人的观测-动作对映射至统一的潜空间，为跨任务、跨平台的技能迁移提供了共享表征。在大规模多模态数据集上对 VLA 基模进行预训练与后训练微调，验证了其在处理复杂语义指令时的泛化性能。针对真机部署中推理频率与执行频率不匹配导致的运动抖动问题，引入了基于三次多项式插值与滑动窗口滤波的后处理平滑机制。实验结果表明，该机制有效抑制了时序预测中的高频噪声，确保了机器人在执行叠衣、抓取等精细化家务任务时的运动连续性与稳定性。

\section*{二、未来展望}
针对本文已有研究的不足，对未来工作提出以下几点展望:

(1) 书法技能模仿学习框架存在以下局限性。框架假设笔画宽度一阶决定因素为毛笔法向按压力度，将书写速度引起的墨水渗透效应视为二阶残差，速度与压力的完全解耦在仅依赖静态图像输入的条件下不可能实现，未来可引入动态书写视频或加装触觉传感器来突破此约束；笔画分割算法基于"书写方向在交叉处不突变"假设，在行书草书多笔画缠绕等极端场景下失效，且仅在艾利特EC66单一平台上测试了有限楷体汉字。未来可在适应更多风格汉字的方向发展该框架。

(2) ConstrainNet框架的设计定位是软约束引导机制，提供特征层面的动作分布偏移而非可证明的硬安全保证。约束观测输入为低维代理表征（2D边界框或关节空间表达），深度方向避障隐式依赖策略网络的视觉泛化能力；对需绝对安全保障的场景，建议结合控制屏障函数等方法。在解决组合约束冲突时当前采取重采样方案，未来可从训练层面在损失函数中加入约束表征空间的正则化项。

(3) VLA系统的后处理机制在高速动态场景下存在适用边界。固定窗口大小（$W=5$）的滑动平均滤波器引入约200ms群延迟，在准静态任务（如叠毛巾）中可接受，但在末端速度超过1.0m/s的动态场景下将引入20cm以上物理偏差。未来可采用自适应窗口策略或BID、RTC等异步推理架构，使动作预测与执行解耦。

(4) VLA跨构型迁移实验目前仅在星海图R1 Lite单一新构型上验证，更大规模跨平台迁移能力仍需检验。潜动作编码器的对比学习损失未显式建模动力学参数，虽是为保持潜空间构型无关性的主动设计，但动力学差异可能对精确力控任务产生潜在影响。未来可引入如力信息以及世界模型等方法填补动力学建模的缺失。